\chapter{绪论}
\label{cha:intro}

\section{研究背景与意义}

脑机接口能够在不依赖外周神经和肌肉输出的条件下建立人脑与外部设备之间的信息通道，在神经康复、辅助交互、智能可穿戴和新型人机协同等方向具有重要应用价值 \cite{chaudhary2016brain, lotte2018review}。其中，基于脑电（Electroencephalography, EEG）的非侵入式脑机接口因采集成本较低、使用风险较小、便于重复实验和系统集成，成为当前研究与应用开发最活跃的技术路线之一。对于希望进入辅助交互与便携式使用场景的系统而言，安全性、可重复性与可集成性往往比单次实验中的极限性能更具现实意义，这也产生了面向真实应用条件构建无创脑机接口系统的研究需求。

\begin{figure}[htbp]
  \centering
  \includegraphics[width=0.92\textwidth]{figures/C1-1_BCI_system_and_apps.png}
  \caption{EEG-BCI系统的基本组成与典型应用场景示意\cite{gu2021eeg}}
  \label{fig:C1-1}
\end{figure}

如图\ref{fig:C1-1}所示，EEG-BCI系统通常由信号采集、预处理与特征提取、解码与控制以及应用终端等环节构成，其应用场景也已由实验室验证逐步扩展到康复辅助、交互控制与状态监测等方向。这种“从可测到可用”的演进进一步凸显了从系统链路角度开展研究的意义。

从任务范式看，脑电脑机接口已经形成了较为清晰的技术谱系。以 P300、稳态视觉诱发电位（Steady-State Visual Evoked Potential, SSVEP）和运动想象（Motor Imagery, MI）为代表的范式，分别对应事件相关模式检测、持续视觉刺激驱动识别和主动意图解码等不同交互路径 \cite{farwell1988talking, bin2011high, pfurtscheller1999event}。其中，运动想象无需真实肢体动作即可诱发可识别的感觉运动节律变化，也不依赖持续外部刺激维持，因此长期被视为主动控制类脑机接口中的典型范式；但与此同时，它对个体训练、状态稳定性和解码策略的依赖也更为突出。

从技术任务看，脑电解码大体围绕事件相关模式检测、脑状态分类识别以及连续意图或认知状态估计展开。面对这些不同的任务或范式，低信噪比和非平稳性始终是脑电解码的主要挑战：一方面脑电信号幅值较弱，易受肌电、眼电和环境噪声污染，整体信噪比较低 \cite{goncharova2003emg, lotte2018review}；另一方面已有研究表明，跨被试差异、跨 session 漂移以及试次级信号质量波动都会直接影响模型泛化与解码稳定性 \cite{lotte2018review, kwon2019subject, li2024joint, Zho16}，表现为脑电信号的非平稳性。

近年来，围绕上述难题，脑电解码研究一方面沿着更深的卷积网络、更复杂的自
注意力结构以及更大规模训练数据持续推进，另一方面也出现了面向跨被试迁移
和表示复用的大规模预训练模型。相关工作试图通过大规模跨数据集学习缓解异
质性、样本不足和泛化困难 \cite{Jia24, Wan24b}。例如，Jiang 等提出的
Large Brain Model for EEG 试图利用大规模跨数据集训练学习更通用的脑电表
征，以提升多任务场景下的迁移能力 \cite{Jia24}；Wang 等提出的 CBraMod
则进一步强调在不同 EEG 格式、不同任务和不同数据来源之间建立统一的表征空
间 \cite{Wan24b}。这类工作说明，大规模预训练和基础模型路线已经开始从单
一任务优化转向更广义的表示复用与跨域泛化。但现有综述也指出，这类路线的
主要优势仍集中于多导联、离线基准和较重适配条件下的表示迁移，其收益并不
能自动等同于低成本、少导联、实时交互和可部署系统中的直接可用性
\cite{Liu26b}。

随着辅助交互终端、可穿戴设备和日常化人机协同需求的发展，无创脑机接口系统正在逐步突破传统高密度有线实验平台的使用边界。一方面，市场上已经出现了以头环、头带和耳周设备为代表的轻量化 EEG 产品，相关产品多聚焦于睡眠监测、专注训练和基础认知调节等场景；另一方面，学术界也在持续探索耳电、前额阵列、低准备度无线采集和针对特定任务的可穿戴 EEG 系统 \cite{He23, Nis22, Arn20, Deb15, Bio22}。这说明“更轻、更易用、更可持续部署”的 EEG 形态并非凭空设想，而是已经在市场探索和学术研究中同时展开的现实趋势，也进一步产生了围绕可穿戴实现、持续运行能力与部署可行性开展系统研究的现实需要。

但便携化、无线化和边缘化并不等同于对实验室设备进行简单的小型化改造。面向真实使用场景，系统往往需要同时兼顾穿戴便利性、布设时间、无线传输、持续运行、本地响应以及维护成本等多重要求 \cite{Cas19c, Nis22, He23}。在这种条件下，系统通常难以继续沿用实验室中的高密度有线配置，而需要围绕佩戴复杂度、部署成本与实时响应等因素重新组织采集与计算链路。相应地，采集端更可能采用简化导联配置和紧凑硬件实现，后端则必须在有限算力、有限存储和有限功耗预算下完成稳定推理。

由此带来的挑战可以概括为两个彼此关联的问题。首先，需要在受限采集与受限资源条件下同时兼顾解码模型的有效性、紧凑性与端侧部署友好性。在采集配置简化的条件下，系统可利用的空间观测通常更为有限，信号质量也更容易受到接触状态、运动伪迹、基线漂移和试次质量波动等因素影响。因此，相关研究不仅需要关注解码是否仍然可行，还需要进一步评估在有限算力和功耗预算下，何种任务目标与模型结构能够满足实际部署要求。同时，当系统进一步走向在线交互时，交互延迟、反馈机制和用户适应会继续改变任务表现，系统关注的重点也不再只是继续提高离线准确率，而是如何在真实使用条件下保持足够稳健、足够及时和足够可持续的交互能力。

本课题以国家科技创新2030--“脑科学与类脑研究”重大项目《侵入式/非侵入式双向闭环脑机接口关键技术及系统》为背景，面向辅助交互与便携式使用场景下的无创脑机接口需求，以运动想象这一典型脑电交互任务为切入点，围绕真实应用中的系统实现与任务落地问题展开研究。结合前述分析，本文重点关注三个相互衔接的问题：其一，面向便携式与可持续运行需求，如何建立可用的双导采集、无线传输与基础交互链路；其二，在受限采集与受限资源条件下，标准高密度离线研究结论应如何收束为更适合真实系统的任务目标与解码方法；其三，离线算法主线如何进一步进入系统集成、模型更新与在线联动过程，并形成具有实际参考意义的初步验证。围绕上述问题，论文分别从双导采集硬件实现、面向低通道端侧场景的解码方法设计以及系统集成与在线联动初步验证三个层面展开工作，旨在建立从标准高密度设置到双导可部署方案的研究路径，为无创脑机接口在真实辅助交互场景中的应用提供参考。

\section{相关研究现状}
\subsection{脑机接口系统与脑电采集研究现状}

从系统实现角度看，脑机接口首先依赖稳定的脑电采集与处理链路。已有综述通常将脑电采集设备划分为导联接口、预处理电路、模数转换电路、处理器、存储与通信模块以及供电模块等部分，其中前端放大与滤波、参考方式、模数转换精度和主控/通信架构共同决定了系统的噪声水平、实时性与可扩展能力 \cite{He20, Martins98, Kim22wearable}。这意味着，对脑机接口系统的研究不仅涉及信号能否被采到，还涉及采集链路能否在有限体积、有限功耗和连续使用条件下稳定工作。

\begin{figure}[htbp]
  \centering
  \includegraphics[width=0.88\linewidth]{figures/C1-2_EEG_acquisition_block_diagram.png}
  \caption{脑电采集设备的一般组成框图（据\cite{He20}整理）}
  \label{fig:c1_2_eeg_acq_block}
\end{figure}

围绕这一链路，现有研究大致沿三条方向展开。首先，传统脑电采集设备研究重点关注前端放大、滤波、参考方式与高精度模数转换等关键环节，以提高微弱脑电信号采集时的共模抑制能力、输入阻抗和噪声控制水平 \cite{He20}。随着高集成模拟前端芯片和便携处理节点的发展，相关工作又进一步转向体积更小、功耗更低、便于无线回传的实现方案 \cite{Martins98, Kim22wearable}。这一路线回答的是“如何把脑电可靠采到”，并为后续可穿戴系统提供硬件基础。

其次，越来越多研究开始探索低准备度与可穿戴化的 EEG 形态，例如头环、头带、耳周 EEG、前额阵列以及面向特定任务的减少导联系统 \cite{Arn20, Deb15, Bio22}。这类工作表明，在睡眠监测、长期状态评估和轻量交互等场景中，更易佩戴、布设更快且适合连续使用的 EEG 设备已经具备明确需求与技术基础。与此相对应，另一类综述和系统研究则从真实应用约束出发，强调 setup burden、穿戴舒适性、无线传输、维护成本、续航和信号稳定性之间的综合权衡 \cite{Cas19c, Nis22, He23}。因此，系统研究的重点正在由单纯提高采集完整性，逐步转向如何以可接受方式反复、长期并稳定地完成测量。

再次，在系统级实现层面，已有研究开始将采集前端、无线链路与边缘算力节点放在同一框架下考察。Wang 等、Schneider 等和 En\'eriz 等分别在 MCU、PULP 和 Arduino 级平台上验证了紧凑型 MI 脑电模型的部署可行性 \cite{Wan20, Sch20, Ene23}；Wan 等通过通道选择与 8 bit 量化结合，展示了在压缩输入与模型规模后仍保持可接受解码表现的可能性 \cite{Wan22b}；部分相邻时序任务也开始围绕 RK3568、NPU 或异构推理引擎开展优化 \cite{You24, Ism23}。这些工作说明，脑机接口系统研究已不再局限于传感器形态本身，而是逐渐延伸到采集、传输、推理和交互节点的整体化协同。

总体来看，现有研究已经从采集链路、可穿戴形态和部署节点三个层面证明了便携式脑机接口的基本可行性，并为后续双导硬件与端侧原型设计提供了重要视野。然而，多数工作仍主要回答某一局部问题，例如前端是否足够低噪声、设备是否更易佩戴，或模型能否在指定节点运行；对于真实辅助交互场景下究竟应采用何种采集配置、何种系统链路以及何种任务组织方式，仍缺少统一而连续的研究路径。

\subsection{脑电解码方法研究现状}

从方法演进看，脑电解码研究已经形成了较为清晰的技术谱系。对于运动想象脑电这一典型任务而言，现有工作既追求更高的离线分类性能，也越来越重视在跨被试、跨会话和真实使用约束下的泛化能力与可用性。若从研究思路划分，当前方法大致可以归纳为四条彼此衔接的路线：传统特征工程与浅层分类、卷积神经网络与紧凑深度模型、自注意力与 Transformer 结构，以及预训练、迁移与在线适配增强。

\textbf{（1）传统特征工程与浅层分类路线。} 早期运动想象脑电解码主要依赖人工特征与浅层分类器。以共空间模式（Common Spatial Pattern, CSP）及其滤波器组扩展方法（Filter Bank CSP, FBCSP）为代表的空间滤波路线，能够突出不同类别在空间方差分布上的差异，并与线性判别分析（LDA）、支持向量机（SVM）等浅层分类器结合取得较高分类效率 \cite{ang2012filter, lotte2018review}。这类方法实现简单、计算代价较低、可解释性较强，至今仍是运动想象任务的重要比较基线。不过，它们通常较依赖预设频带、人工特征和相对稳定的采集条件，因此在跨被试迁移、复杂噪声场景与长期在线使用条件下存在明显边界。

\textbf{（2）卷积神经网络与紧凑深度模型路线。} 为减少对手工特征工程的依赖，卷积神经网络逐步成为脑电解码的重要技术路线。Schirrmeister 等提出的深浅卷积网络说明，端到端学习方式能够直接从原始脑电片段中提取具有判别性的时空特征 \cite{schirrmeister2017deep}；Lawhern 等提出的 EEGNet 进一步通过紧凑型卷积结构在较低参数规模下兼顾跨范式适应能力与部署友好性 \cite{lawhern2018eegnet}。此后，多分支卷积、时频联合卷积等变体持续出现，其共同特点是在保留局部时空建模能力的同时，尽量控制参数量与推理复杂度。这一路线推动脑电解码由人工特征设计转向自动表示学习，也为后续端侧友好建模提供了重要基础。

\textbf{（3）自注意力与 Transformer 路线。} 在 CNN 基础上，Transformer 及其变体被进一步引入脑电解码领域，以增强对长程时序依赖和跨特征关联的建模能力 \cite{vaswani2017attention}。Song 等提出的 EEG Conformer 通过卷积与自注意力协同提取局部时空特征和全局时间依赖 \cite{songEEG2022}，Xie 等展示了 Transformer 结构在原始脑电分类任务中的应用潜力 \cite{xie_transformer-based_2022}。随后，研究重点又从“能否使用自注意力”转向“如何在控制复杂度的同时保留关键建模能力”。例如，Wan 等将通道注意力与窗口化机制结合建模时空耦合信息 \cite{wan2023novel}，Hua 等通过卷积与滑动窗口注意力的协同设计增强多尺度特征建模能力 \cite{hua2023improved}，Keutayeva 等则进一步强调了紧凑型结构在跨被试任务中的部署潜力 \cite{keu2024compact}。这表明，自注意力路线已从单纯追求更强表达能力，逐步转向更加关注高效性、紧凑性与可部署性。

\textbf{（4）预训练、迁移与在线适配增强路线。} 近年来，围绕异质性、样本不足和泛化困难，预训练与大规模表示学习开始成为脑电解码研究的新方向。相关工作试图通过大规模跨数据集学习缓解跨被试差异、标注稀缺和表示复用问题 \cite{Jia24, Wan24b}。例如，Jiang 等提出的 Large Brain Model for EEG 试图利用大规模跨数据集训练学习更通用的脑电表征，以提升多任务场景下的迁移能力 \cite{Jia24}；Wang 等提出的 CBraMod 则强调在不同 EEG 格式、不同任务和不同数据来源之间建立统一的表示空间 \cite{Wan24b}。与此同时，在线研究持续指出，仅靠更强的离线表征并不足以直接解决实际交互问题。Pfurtscheller 等和 Neuper 等的早期研究已经表明，反馈训练会显著影响左右手感觉运动节律的可分辨性 \cite{Pfu97, Neu99}；Meng 等进一步将导联数量、算法结构和长期训练表现放在统一框架下考察，说明紧凑通道方案在合理任务设计下仍可获得具有竞争力的在线表现 \cite{Men18}。因此，解码方法的有效性并不完全由离线模型结构决定，表示迁移、用户训练、校准代价和系统链路都会改变最终表现。现有综述也指出，预训练与大模型路线的主要优势仍集中于多导联、离线基准和较重适配条件下的表示迁移，其收益并不能自动等同于受限采集和在线交互系统中的直接可用性 \cite{Liu26b}。

总体来看，脑电解码方法已经形成了从传统空间滤波、卷积建模到自注意力结构，再到预训练与在线适配增强的清晰演进路径，并在标准数据集上不断提升分类性能。然而，现有证据仍主要建立在离线、标准协议和相对高密度导联条件之上，关于这些方法在真实系统约束下如何完成任务收束、结构压缩与在线协同，仍缺少更统一的回答。这也正是后续不足分析与本文研究路线需要进一步展开的部分。

\section{现有研究存在的不足}

综合前述研究可以看出，当前运动想象脑电解码虽然在模型设计、局部高效注意力和轻量化探索方面取得了明显进展，但对于辅助交互与便携式使用场景下的低成本真实系统目标而言，仍存在以下四方面不足。

第一，场景驱动的系统约束仍缺少整体化分析。现有研究往往分别讨论低成本采集、边缘部署或在线交互中的某一方面，却较少从辅助交互与便携式使用场景出发，统一分析穿戴便利性、无线通信、持续运行、本地响应以及采集稳定性之间的耦合关系。结果是系统需求与算法目标之间的对应关系仍不够清晰。

第二，现有标准任务与真实少导联系统目标之间仍存在明显错位。公开数据集中的高密度四分类设置为算法研究提供了统一基准，但当导联数显著降低时，可观测信息边界随之收缩，原有任务定义并不一定仍然合理。许多工作默认“模型只要足够强就能迁移到更少通道”，却较少系统回答在双导或少导联条件下究竟应保留什么目标、舍弃什么目标，以及应如何构建相应的定量证据链。

第三，算法设计与端侧部署约束之间的协同仍然不足。现有研究往往分别强调高离线精度、模型压缩或硬件加速中的某一方面，但较少在统一框架下同时考虑输入通道规模、结构表达能力、参数规模、量化友好性和推理延迟等约束 \cite{keu2024compact, wan2023novel, hua2023improved}。尤其是在运动想象脑电场景中，“能在公开数据集上取得较高准确率”并不等同于“能够在少导联端侧系统中稳定运行并保持可用表现”。

第四，在线闭环证据与持续校准机制仍有待进一步明确。便携式脑机接口不仅取决于模型推理性能，还同时受到采集链路噪声、BLE 传输时序、试次间漂移和在线校准代价等因素影响。已有研究说明反馈训练、个体差异和试次质量会直接影响在线表现，但关于如何在低成本设备上通过够用的模型、轻量的在线更新和合理的交互设计获得稳定收益，公开研究仍然有限。这也是本文后续在系统集成与在线初步验证部分进一步关注的问题。

\section{本文主要研究内容与技术路线}
\subsection{主要研究内容}

围绕上述问题，本文面向辅助交互与便携式使用场景下的低成本无创脑机接口系统，开展了以下几方面研究工作。

（1）设计并实现双导脑电采集硬件系统。针对辅助交互与便携式使用场景下的低成本采集需求，本文完成了以双导联采集为核心的硬件系统设计与实现，构建了脑电前端、主控、BLE 回传和上位机交互的基础链路，并通过节律验证、标签链路验证和通信稳定性测试，为后续在线实验提供系统基础。

（2）研究面向低通道端侧场景的改进型自注意力脑电解码方法。针对标准高密度解码模型难以直接迁移到少导联端侧系统的问题，本文设计了兼顾局部时序建模能力与部署友好性的改进型自注意力解码方法，在标准条件下完成性能、复杂度与量化友好性分析，并为后续低通道扩展提供算法基础。

（3）构建从标准高密度设置到双导可落地任务的收束路径。围绕 `22导 / 4分类` 标准设置，本文系统分析了导联压缩下的性能退化边界，并在此基础上将目标进一步收束到 `2导 / 2分类` 的左右手运动想象解码任务，从而使算法目标与双导系统可观测信息边界相匹配。

（4）完成双导端侧系统集成与在线联动初步验证。本文以双导硬件为前端、以上位机和可扩展推理节点为原型载体，构建了预训练初始化、轮次级模型更新与跨轮次判别反馈的在线联动机制，并结合伪在线评估和真实在线试跑，对系统的可实现性进行了初步验证。需要指出的是，当前实验阶段的在线验证载体主要为 Linux PC 原型环境，后续目标部署节点则面向 RK3568 / RKNN 一类嵌入式平台。

\subsection{技术路线与论文组织}

全文的技术路线与组织结构如图\ref{fig:org}所示。
\begin{figure}[htbp]
  \centering
  % \includegraphics[width=0.8\textwidth]{org_structure}
  \caption{论文技术路线与组织结构}
  \label{fig:org}
\end{figure}

本文后续各章安排如下：

第 2 章介绍运动想象脑电解码与端侧系统实现所需的关键技术基础，并在比较可选技术路线的基础上说明本文采用的建模与部署路线。

第 3 章围绕双导脑电采集硬件系统的设计与实现展开，重点说明采集前端、主控、无线通信和基础链路验证结果。

第 4 章研究面向低通道端侧场景的改进型自注意力脑电解码方法，在标准条件下建立离线基线，并进一步分析导联压缩下的性能退化边界，完成从 `22导 / 4分类` 到 `2导 / 2分类` 的任务收束与模型适配。

第 5 章进一步将第 4 章得到的低通道解码主线映射到双导系统中，介绍系统集成、在线实现机制、伪在线评估以及真实在线初步验证结果。

第 6 章对全文工作进行总结，并围绕更强的人机协同、自适应学习和更自然的实验范式给出后续展望。
